% AI-written, AI cross-reviewed; no human sign-off; not a claim that RH is proved
% RH Notes blueprint source excerpts for Z4.3
% Verbatim excerpts in the notation of the parent blueprint. The labelled editions and source ranges appear below.
% See the unit YAML for the current proof status. Earlier-round passages retain their original wording.

% Start of atlas edition annotation. The following block is outside the original excerpt.
\paragraph{Status of named inputs in this reading.}
The first status refers to the input of the common-good-model construction. Hensel lifting is used in the later local calculation. The reference to a Weil operator names an object in that calculation.
\noindent\nolinkurl{GlobalTheory.Published.GlobalArtinReciprocity} --- hypothesis at this head (\path{Development/AbelianVarietyTheory/PublishedTheorems/AutomorphicForms₂/ExistsUniqueHcharOfFiniteOrder.lean:70--72}); not compared.\par
\noindent\nolinkurl{GlobalTheory.Corev.CmTwoClosing.cmGoodModel_of_globalArtinReciprocity} --- in-repo constant (\path{Development/AbelianVarietyTheory/Proofs/DecoratedLatticeProofs/CmGoodModelOfGlobalArtinReciprocity.lean:10--11}); not compared.\par
\noindent\nolinkurl{GlobalTheory.T.chebotarev_degreeOne_valuationSubrings} --- in-repo constant (\path{Development/AbelianVarietyTheory/Proofs/DecoratedLatticeProofs/ChebotarevDegreeOneValuationSubringsStatement.lean:5--6}); not compared.\par
External theorem numbers in the original excerpt are not verified here unless a separate edition-specific source record documents the comparison.
% End of atlas edition annotation.

\section*{Z4.3. Off-level kernel valuations}

% Parent source SHA-256: 1af64922b98ed2310744411b1f18b25cc426c97220646321471836e1a3a409d6
% blueprint/ch_P.tex lines 12700 through 12975; r6:Z4.3
\subsection*{Z4.3. Rational class pairs, complex forms and the common CM model}
\phantomsection\label{r6:Z4.3}
Fix the true columns, the standard tie for the selected kernel $t$ on
$\mathrm{ex}$, and the twist identification of
Paragraph~\ref{par:r5-z43-ties}. Write $F$ for the coefficient field,
$\sigma:F\hookrightarrow\mathbf C$ for its given embedding and $C$ for
the twisted chain of a word $u$ and a used tuple $\mu$. Two statements
have different quantifiers. For each off-level place, each $i<m$, each
pair $c,c'$ and each rational source vector, one needs a rational
target vector for the entries of $t$. For the complex form spaces one
needs the operator of $C$ to carry $F_i$ into $F_{i+1}$, both on the
used twists and on the member chain.
See \path{KmapRat}, \path{KFi} (\path{Formalism/SharpRestriction/Divisibility/OrbitPlaces.lean:107--109,1367--1369})
and \path{at_of_true_columns} (\path{Development/AbelianVarietyTheory/Proofs/SatakeParameters/AtOfTrueColumns.lean:30--47}).

\paragraph{The rational vector for an individual pair.}
Let $\nu_i(\tau,r)$ be the integral polynomial weight of the class
kernel and $k_i(\tau)$ its last coordinate, for $i<m$.
Antitonicity gives $0\le k_i(\tau)\le\nu_i(\tau,r)$.
Use the weight with natural part $\nu_i-k_i$ and evaluation shift
$s_i$, where $s_0=k_0-s^\sharp_0$ and $s_{i+1}=-s^{\rm div}_i$.
At an interior step the old weight coordinates increase by
$s^\sharp_{i+1}-s^{\rm div}_i$ and the new coordinates all equal
this increment. In fact the indices $N_i+a$, $0\le a<t_0+1$,
have quotient $i$ and remainder $a$ after subtracting the initial
rank and dividing by the rank increment $t_0+1$. Consequently the
natural part at $i+1$ is precisely $\nu_i$ padded by zeros,
and $s_i+s^\sharp_i=k_i$. At the terminal step this padding is
the definition of the last weight. Thus this calculation includes
$i=m-1$. See \path{nu_step_old}, \path{nu_step_new}, \path{wtTie_meet},
\path{natW_succ} (\path{Development/AbelianVarietyTheory/Definitions/GaloisGeometry/KernelRationalReferenceTuple.lean:17--63,65--132,133--215,346--420}).

Here is the evaluation comparison over $F$. Write
$D_k(Y)=\prod_\tau\det(Y_\tau)^{k(\tau)}$. In each block the
bideterminants spanning the dual Weyl space have exactly as many
full-size factors as the last coordinate of their shape. A full-size
factor either repeats a row and vanishes or equals the determinant
up to sign. Hence the indicated power divides every generator.
Taking products over blocks and finite linear combinations gives
$D_k\mid Q$ for the weight-$\nu$ polynomial $Q$. Write $Q=D_kP$.
Under invertible triangular column substitution, $D_k$ acquires
the character $\prod_\tau(\prod_r\beta_{\tau,rr})^{k(\tau)}$.
This scalar and $D_k$ are nonzero (evaluate the latter at identities).
Cancelling them proves the required weight $\nu-k$ for $P$.
On the initial frame $\det^{-s^\sharp_0}Q$ thus evaluates as
$\det^{s_0}P$ block by block, and multiplication by $D_k$ gives
the converse inclusion. At higher levels the preceding padding
identifies the polynomial spaces directly. Quotienting by the kernel
of evaluation preserves the range since every coset has a
representative. This identifies the range of the \emph{actual} entry
evaluation at both ends by the equality in $\mathrm{IsEntC}(t)$.
See \path{range_evC}, \path{vqTie_kerEntC},
\path{gdet_pow_dvd_dualWeyl}, \path{dk_dvd_wtPol},
\path{frameWeightPolyDivByDet_proof}
(\path{Development/AbelianVarietyTheory/Definitions/GaloisGeometry/KernelRationalReferenceTuple.lean:442--531,545--608,624--674}).

Fix $v$ and choose $P$ representing its shifted evaluation. Put
$z=\rho_i(c)\rho'_i(c')$. The two factors have norm one.
For a derivative index $a$ of multidegree
$d_\tau=\sum_r\nu_i(\tau,r)$ in block $\tau$, choose $J_a\in F$
with $\sigma(J_a)=z\,\mathrm{jet}_a(\Theta_{cc'})$. To see why the
specified class kernel has this property, use the support and
period of its finite datum to sum over the finite quotient of the
divided hom lattice by the multiplied hom lattice. The individual
translates are smooth, so the jet of the sum is the sum of the jets.
The determinant-character identity and the weights of the finite
datum put the remaining scalar factors in $\sigma(F)$. For each
representative the true-column hypothesis \path{JetsIn} gives the
required field element for the remaining product at degree $d_\tau$.
The equality between that degree and the true-column degree is obtained
by taking natural floors of the nonnegative integral weight entries
and summing. The standard tie identifies the differential generators
and the theta function of $t$ with these true-column data.
See \path{rho_jet_mem}, \path{rk_of_std}
(\path{Development/AbelianVarietyTheory/Proofs/SatakeParameters/TrueKernelJetsMulDetCharactersMemCoefficientFiel.lean:24--68},
\path{Development/AbelianVarietyTheory/Proofs/SatakeParameters/CoefficientFieldConditionOfStandardKernelCompari.lean:18--55}).

Set $J_a=0$ at other degrees. The differential polynomial formed
from $D_{k_i}P$ only uses multidegree $(d_\tau)_\tau$. Thus its jet
map is unchanged after applying $\sigma$ to this choice of
coefficients. Unfolding the pair operator on the polynomial model
and using $z\ne0$ identifies its complex output with the evaluation
of $\mathrm{jetMap}(J,k_i,P)$, with the target shift
$-s_i^{\rm div}$. The target weight is the padded weight just
computed. The range equality therefore supplies $v'$ in the
rational target quotient with exactly the required evaluation.
The argument works for each off-level place in the definition.
See \path{wtTie_meet}, \path{kmapRat_of_vqU},
\path{kmapRat_kerEntC'} (\path{Development/AbelianVarietyTheory/Definitions/GaloisGeometry/KernelRationalReferenceTuple.lean:133--215,244--278,685--687}).

\paragraph{The stabilizer calculation on the same chain.}
Use the true-column standard reader in the auxiliary tied record
$t'$. It is constructed with $I.\mathrm{kerRd}$ and its standard
top chain; the lattice and theta comparison equations are transported
to this reader. It is a record on the \emph{same} $\mathrm{ex}$.
In its frame the complex Young span is the polynomial value space.
The action fixing a class is given by left multiplication by the
matrix of its stabilizer with the attached character. Thus the two
tests for membership at a positive level are the polynomial value
test and this transformation law, with the actual frame and weight
of the true columns. The standard level reader identifies its
stabilizers with these specified points. See \path{stdKer_stdRd},
\path{stab_pin}, \path{formsAre_stabStd}, \path{duty_of_std}
(\path{Development/AbelianVarietyTheory/Proofs/SatakeParameters/StandardKernelDataWithStandardBaseLevels.lean:14--36},
\path{Development/AbelianVarietyTheory/Proofs/SatakeParameters/TrueChainStabilisersEqStandardBaseLevelPoints.lean:24--65},
\path{Development/AbelianVarietyTheory/Proofs/SatakeParameters/TrueChainComplexFormsStabiliserLaw.lean:18--57,69--117},
\path{Development/AbelianVarietyTheory/Proofs/SatakeParameters/ConditionsAtStandardKernelComparisonTrueColumns.lean:16--24}).

Fix a target stabilizer represented by $\gamma$ with multiplier $k$
over the finite coefficient field. It preserves the target lattice.
The test function $\Xi$ is invariant under addition of the level
multiple of the hom lattice. Before the first coupled step use the
source identity and change only the target frame. At the first
coupled step the stabilizer additionally fixes the coupling class,
so its chosen representative can change by an allowed period. At the
last step apply the top-data law at a level prime and the double-depth
congruence at the other primes; the inverse target-frame matrix is
the level-group element there. Multiplying the local identities gives
$\Xi(\gamma x)M=\sigma_0(k)\Xi(x)$, with $M$ the product of the
embedded finite character values at $\det\gamma$. These are the
three cases $i<m_2$, $i=m_2$ and $i=m_2+1$.
See \path{term_law}, \path{xiEntry_stab}, \path{xiCpl1_stab}
and the top-level product calculation
(\path{Development/AbelianVarietyTheory/Proofs/SatakeParameters/ClassSumTermLawStabiliserTarget.lean:25--70}, \path{Development/AbelianVarietyTheory/Proofs/SatakeParameters/XiEntryStab.lean:13--37},
\path{Development/AbelianVarietyTheory/Proofs/SatakeParameters/XiCpl1Stab.lean:21--54}, \path{Development/AbelianVarietyTheory/Proofs/SatakeParameters/LastStepDatumLawTopLevelStabilizer.lean:87--156}).

Write $d=\det\gamma$. It is a norm-one unit. The global character
identity gives $M=\prod_\tau\tau(d)^{-e_i(\tau)}$ with
$e_i=s^{\rm div}_i+N_i/2$. The common half-rank factor drops out
because $\prod_\tau\tau(d)=1$. In particular $M\ne0$.
Reindexing the finite lattice sum by $\gamma$ uses the first covariance law
at a target stabilizer (Section~\ref{r14:equivariance}). Under its suppliers
this makes the entire theta function
an eigenfunction of left multiplication with eigenvalue
$\sigma_0(k)/M$. Its differentiability follows by taking the finite
sum of the differentiable translated vacua and composing with the
linear change of archimedean coordinates. Applying the polynomial
differential operator at zero after changing the target frame
multiplies the answer by this same eigenvalue, by the chain rule
with the inverse target matrix. The target determinant in the
denominator contributes precisely $M$. Hence the pair term acquires
$\sigma_0(k)$, including its character. Its polynomial value has the
target weight by the same padded-weight jet calculation over
$\mathbf C$. See \path{termLaw_of_xi}, \path{thC_left_stab},
\path{kerTie_entire}, \path{applyAt0_polD_eigen}, \path{lawAt_Kmap}
(\path{Development/AbelianVarietyTheory/Proofs/SatakeParameters/PairKernelLawOfFiniteDatumLaw.lean:21--59,85--121}, \path{Development/AbelianVarietyTheory/Proofs/SatakeParameters/ThCLeftStab.lean:368--420},
\path{Development/AbelianVarietyTheory/Proofs/SatakeParameters/KerTieEntire.lean:12--24}, \path{Development/AbelianVarietyTheory/Proofs/SatakeParameters/PairKernelTermTargetStabiliserLaw.lean:18--75},
\path{Development/AbelianVarietyTheory/Proofs/SatakeParameters/PairKernelMapsStepWeightValues.lean:19--33}).

For each $c'$ sum the pair terms with the source weights $C.\mathrm{wt}_i(c)$.
Both tests are linear, proving membership in $F_{i+1}$.
At the reference tuple the given isomorphism of rational lattices
intertwines their evaluations on the common chain. The spans of all
evaluated vectors agree at $i$ and $i+1$ (expand in either lattice
basis), giving the member-chain assertion as well. This proves
$\mathrm{KFi}(D,\mathrm{ex.tw})$ from $t'$, while the rational class-pair
assertion above is on $t.\mathrm{ent}$. These are precisely the two
separate arguments required by the rational kernel form.
See \path{K_mem_Fi}, \path{carries_one}, \path{kFi_of_tie},
\path{KernelFormRational} (\path{Development/AbelianVarietyTheory/Proofs/SatakeParameters/KernelOperatorMapsStabiliserEquivariantForms.lean:11--38},
\path{Development/AbelianVarietyTheory/Proofs/SatakeParameters/MemberFormsStableOfReferenceTuple.lean:7--47}, \path{Development/AbelianVarietyTheory/Proofs/SatakeParameters/KernelClassSumsCarryChainComplexForms.lean:13--25},
\path{Formalism/SharpRestriction/Divisibility/OrbitPlaces.lean:1374--1375}).

\paragraph{The common good model and its uses.}
The row construction takes global Artin reciprocity as its stated input.
The induced ideal-class map kills inertia and sends arithmetic
Frobenius to the class of the prime. To see where this gives the
unramified CM invariant, take an integral monic polynomial of the CM
value $j$ and remove the finite primes dividing an integer separating
its distinct roots upon reduction. Intersect the open kernel of the
class map with the subgroup fixing all those roots. Degree-one
Chebotarev in this finite quotient represents the coset of a given
inertia element by a Frobenius outside the discarded set. Its ideal
class is trivial. The CM Frobenius congruence there and the arithmetic
Frobenius congruence imply that the root and its transform have equal
reduction. Root separation makes them equal, and equality of the
cosets gives the same conclusion for the original inertia element.
Thus $j$ is integral and its image is in the completed maximal
unramified extension at the required $p$. See \path{inertia_fixes},
\path{good_model}, \path{cmGoodModel_of_globalArtinReciprocity}
(\path{Development/AbelianVarietyTheory/Proofs/SatakeParameters/CmLatticeGoodModelUnramified.lean:10--69}, \path{Development/AbelianVarietyTheory/Proofs/DecoratedLatticeProofs/CmGoodModelOfGlobalArtinReciprocity.lean:10--11}).

There is a further valuation step from $j$ to a short model.
Write $J=\iota(j)$. For $J=0$ and $J=1728$ use $(g_2,g_3)=(0,1)$
and $(1,0)$ respectively, scaling by a sixth or fourth root in
$\mathbf C_p$. In the other cases, normalize a homothetic lattice
so that both invariants equal $C=27J/(J-1728)$. When $|J|<1$ put
$|J|=p^n$; after scaling to discriminant one by $b$ this gives
$|b|^{12}p^{2n}=1$. Take one point from each sign pair in an odd CM
endomorphism kernel of order coprime to $3p$ so its resulting size
$t$ is nonzero modulo 3. In the ramified-at-three case the construction
takes the intersection of the kernels of an auxiliary prime and
$r+\sqrt{-d}$ with prime valuation exactly one in its norm.
The root polynomial of the kernel, or the common roots of the two
kernel polynomials, puts the product of their $\wp$-coordinates in
the unramified field. Its norm is $p^a$ for integral $a$. Each scaled
coordinate is a unit by the prime-to-order torsion test at the
exceptional reduction, so $|b|^{2t}p^a=1$. Eliminating $|b|$ gives
$nt=3a$, hence $3\mid n$. When $|J-1728|=p^n<1$ the discriminant
calculation instead gives $|b|^4=p^n$. The same argument with $t$
odd gives $nt=-2a$ and $2\mid n$. In the unit cases the respective
exponent is zero. Write $|J|=p^{3m}$, $|J-1728|=p^{2n}$ and put
$h=p^{m-n}$. The model $(h^2 C,h^3 C)$ has discriminant
$h^6 C^2(C-27)$ of norm one; the cubes and squares of its
coefficient norms are $|J|$ and $|J-1728|$. The common scaling exists
because $C$ is the ratio of the cube and square of the original two
invariants. This gives the witnesses for CM2 at $p\ge5$, $p\nmid d$,
for the specified $\iota$. See \path{kerProd_prin}, \path{kerProd_ideal},
\path{unit_prod}, \path{valDiv3}, \path{valDiv2}, \path{model_generic},
\path{model_of_valDiv} (\path{ReciprocityTheory/ReductiveGroups/UnramifiedIntegrals/RedKerPrin.lean:30--91}, \path{ReciprocityTheory/ReductiveGroups/UnramifiedIntegrals/RedKerIdeal.lean:33--79},
\path{ReciprocityTheory/ReductiveGroups/UnramifiedIntegrals/RedValB.lean:40--52}, \path{ReciprocityTheory/ReductiveGroups/UnramifiedIntegrals/NormWeierstrassCoeffSmallJ.lean:98--185},
\path{ReciprocityTheory/ReductiveGroups/UnramifiedIntegrals/Norm1728EqOneOfFiveLe.lean:45--148}, \path{ReciprocityTheory/ReductiveGroups/UnramifiedIntegrals/RedModel.lean:42--126,128--189}).

For the curve dictionary the additional input is CM4: prime-to-$p$
nonzero torsion gives a nonidentity point on the scaled short curve.
Inertia acts trivially on this prime-to-$p$ torsion in the good
reduction argument. Reduction of a nonidentity such point is
nonidentity, since the formal group has no nonzero prime-to-$p$
torsion. Its two affine coordinates are consequently integral and
unramified. The second coordinate in the short model uses
$b^3\wp'/2$; multiplying by $2$ preserves both conclusions. The
dictionary uses CM2 and this torsion statement; its CM3 parameter is
unused. Conversely the lattice model removes the inert-prime ideal
factor preserving the scaled lattice, then constructs the model of
the coprime ideal from CM2 and CM3; its CM4 parameter is unused.
This dependence is visible in the proof terms, although both public
signatures accept all three propositions. See \path{cmTorsion_of_point},
\path{cm4}, \path{datum_exists_lattice}, \path{curveDict_of_cm}
(\path{ReciprocityTheory/ReductiveGroups/UnramifiedIntegrals/TorStatus.lean:15--38}, \path{Development/AbelianVarietyTheory/Proofs/SatakeParameters/ExistsRingEquivTranscendental.lean:277--351},
\path{Development/AbelianVarietyTheory/Proofs/SatakeParameters/CurveDictOfCm.lean:17--20}).

\emph{The hand proof of the following height-two implication for the specified CM3 model is on page~\pageref{r9:Z4.3}.}
Let $U$ be the completed maximal unramified field at this inert
prime, $\varphi$ its isometric Frobenius lift and $\lambda$ the
normalized short-curve logarithm. The integral two-variable law
$\Phi$ has logarithm $\lambda$, and the integral endomorphism $A$
has zero constant and satisfies $\lambda(A)=a\lambda$,
$|\varphi(a)-a|=1$, $\varphi^2(a)=a$. Its multiplier is the
embedded CM generator. Let $e=[p]$ have zero constant, integral
coefficients, $\lambda(e)=p\lambda$, and a coefficient of norm one
at some degree $i\le p^2$. The comparison for this logarithm is the coefficient estimate
\[
 \exists c\in U,\quad |c|=1,\qquad
 \lambda(X)-\frac{c}{p}(\varphi^2\lambda)(X^{p^2})
       \in\mathcal O_U[[X]] .
\]
Frobenius acts on coefficients here. The source uses the comparison
of the short multiplication-by-$p$ series and its logarithm, then
twisted division to eliminate higher Frobenius iterates in this
estimate. This is the particular implication \path{hon_main},
specialized in \path{cm3}, rather than an unspecified choice of CM
curve. Given its unit $c$, solve $\epsilon=c\varphi^2(\epsilon)$
among units by successive lifts from the algebraically closed residue
field. The formal inverse-logarithm comparison then gives integral
coefficients of
$u_{g_2,g_3}(\epsilon\sum_{j\ge0}X^{p^{2j}}/p^j)$ in $U$.
This is the required integral height-two coordinate, with unit linear
term. See \path{HD}, \path{hon_main}, \path{lang_semilinear},
\path{fin_integral}, \path{cm3}
(\path{ReciprocityTheory/ReductiveGroups/UnramifiedIntegrals/FgDefs.lean:14--28}, \path{ReciprocityTheory/ReductiveGroups/UnramifiedIntegrals/FgHonH.lean:220--256}, \path{ReciprocityTheory/ReductiveGroups/UnramifiedIntegrals/SemilinNormPowSubPow.lean:76--101},
\path{ReciprocityTheory/ReductiveGroups/UnramifiedIntegrals/FgFinA.lean:115--132}, \path{ReciprocityTheory/ReductiveGroups/UnramifiedIntegrals/FgClose.lean:25--91}).

The algebraicity of the translated Taylor coefficients is supplied
separately: a torsion point lies in the saturated lattice, the
normalized translated vacuum has algebraic jets there, and multiplying
by the algebraic character factor preserves algebraicity. In the row
service the one CM2 proposition supplies the three-output theorem,
inert integrality, inert admissibility and the formal-curve existence
argument. With the height-two coordinate, the common curve, basis
and input quantifiers are precisely those of
Paragraph~\ref{par:r4-z43-host}. The place and embedding are fixed before
choosing the local frame, which is then available for every allowed
translate. For the selected row the shared rational unitary covariance, including its target-stabilizer use, still depends on its suppliers. The finite expansion has been hand-derived under the translated-jet and reader hypotheses (the two original questions are on p.~\pageref{r9:z43-open}). The suppliers at the individual callers are registered on p.~\pageref{r10:z43-register}. These hypotheses identify the data to which the class-pair map, complex-form preservation and four-way numerical comparison apply. See \path{algTAll}, \path{servH_row_of_artin},
\path{torsionTaylorAlgebraic_latticeExists_of_cm}, \path{at_of_true_columns}
(\path{Development/AbelianVarietyTheory/Definitions/GaloisGeometry/AlgTAll.lean:12--25}, \path{Development/AbelianVarietyTheory/Proofs/SatakeParameters/BundleBothPartsOfStandardKernelOfGlobalArtinReci.lean:15--24},
\path{Development/AbelianVarietyTheory/Proofs/SatakeParameters/TorsionTaylorAlgebraicLatticeExistsOfCm.lean:15--19}, \path{Development/AbelianVarietyTheory/Proofs/SatakeParameters/AtOfTrueColumns.lean:30--47}).

% Parent source SHA-256: 99a8461100a8be8176578b515f15d452958d95ac91f86554c8727ce0a98ee36e
% blueprint/ch_P_r14.tex lines 1333 through 1509; r14:comb-tied, r14:comb-seam, r14:comb-pseries
\subsection{The selected reader and its $p$-adic combination (Z4.3(ii))}
\label{r14:comb-tied}

Fix a tied entry $t$ on the data $(D,L,h_T,G_i,G_l,K,Y,c_X,\mathrm{ex})$.
Let $u,\mu$ be its word and tuple, $i<m$ and $(c,c')$ any classes
in steps $i,i+1$. Write $r=t.\mathrm{rd}(u,\mu)$, $s=t.\mathrm{seam}$,
$A=t.\mathrm{amb}$, $T=t.B.\mathrm{tor}$ and
$e=t.B.e^{\rm nat}(t.\chi)$. Set $M=r.N_i(c)$, $M'=r.N_{i+1}(c')$,
$L_0=s.\mathrm{lat}(M)$, $L_1=s.\mathrm{lat}(M')$ and
$S=(t.O(u,\mu)).S$. At a level-free place $w$, put
$p=\operatorname{resChar}(w)$ and $\lambda=t.\mathrm{ent}.\mathrm{lam}(w)$.
The assertion to submit is the precise implication of
\nolinkurl{GlobalTheory.Ggeo.KerComb.comb_of_tied}. Assuming
\nolinkurl{AlgTAll} and $t.\mathrm{Hyp}$, the ordered lattice pair is
admitted and the $p$-series of this very class-pair kernel satisfies
\nolinkurl{TaylorLogMetric.IsComb} for $T,e,\lambda,L_0,L_1$.
The site theorem used for Z4.2.04 has a different conclusion, $J.\mathrm{Dict}
\Rightarrow S_{\rm site}.\mathrm{Pair}$. The common calculation is
the complex finite expansion of Section~\ref{r14:pair-kernel}.
See \path{SatakeParameters/CombOfTied.lean:14--33} and
\path{SatakeParameters/KernelSiteLatticeTuplePairsAdmissible.lean:23--45}.

\paragraph{Transport of this kernel to the chosen period lattice.}
\phantomsection\label{r14:comb-seam}
Specialize $t.\mathrm{Hyp}$ at $u,\mu,i$. It gives
$A.\mathrm{sc}=\operatorname{area}(T)^{-1}$ and a positive level
$\ell=r.\mathrm{lev}$ supported on $S$. On the specified adjacent pair
it gives the pushed vacuum equality, the inclusion of pushed periods in
$\mathrm{La}(D,T,L_0,L_1)$, integrality of both forms on the period
lattice, the three nonzero relative indices and $\mathrm{IntOff}(S,r.\Xi(x))$
for every $x$. Here $\mathrm{IntOff}$ means algebraic with norm at most one
under every $p'$-adic embedding off $S$. Apply the volume-one expansion
of Section~\ref{r14:kernel-finite} with these data. In its notation
$\Lambda=\mathrm{LamC}$, $A_-=\mathrm{AC}$, $B_+=\mathrm{BC}$, the
three quotients, in order, are
\[
 \frac{A_-}{A_-\cap\Lambda},\qquad
 \frac{\Lambda\cap A_-}{\Lambda\cap A_-\cap B_+},\qquad
 \frac{B_+}{B_+\cap\Lambda}.
\]
The middle cardinal is the Fourier denominator. The index pair $(x,y)$
ranges over representatives of the first and third quotients. Its weight
$\rho_{x,y}$ is integral off $S$, and the shift $x+y$ has
$\ell(x+y)\in\Lambda$, since both $\ell x$ and $\ell y$ belong to
$\Lambda$. The datum used here is $\mathrm{onV}(r.\Xi)$, which at a
point in the embedding image takes a value of $r.\Xi$ and elsewhere
is zero. These are the inputs of \nolinkurl{theta_coset_kernel};
see \path{FirstKernelLawReading.lean:92--94,114--123,447--448},
\path{XiIndCClass.lean:210--211,224--228,239} and
\path{PeriodLatticeFinitelyGenerated.lean:248--290}.

Write $s_*F(z)=F((z_{s.\mathrm{iT}(\tau)})_\tau)$ and
$P_s v=(v_{s.\mathrm{iT}^{-1}(\tau)})_\tau$.
Reindex the finite sum of matrix entries over $\tau$ in the Hermitian
form. Its coefficient is $\operatorname{area}(T)^{-1}$ in both coordinates,
so the values on $(z,P_sv)$ and $(P_sv,P_sv)$ are the two inner products
in the old exponential multiplier. Reindexing $z+P_sv$ by $s.\mathrm{iT}$ gives the reindexed $z$ plus $v$. Substitution in the formula
for $U$ therefore gives $s_*(U_vF)=U_{P_sv}(s_*F)$. Reindexing commutes
with addition, hence $\ell P_s(x+y)$ lies in $\mathrm{La}(D,T,L_0,L_1)$
by the period inclusion. Applying $s_*$ term by term and the pushed
vacuum equality gives the expansion of $r.\mathrm{Th}$ with the
same weights, shifts $v_{x,y}=P_s(x+y)$, and vacuum $e(L_0,L_1)$.
The defining equalities $t.O_{\rm lat}$ and $t.O_{\rm Th}$ identify
both lattices and the kernel with those of the actual entry.
This proves the reader implication \nolinkurl{reading_is_kernel};
\nolinkurl{kernel_comparison_holds} simply unpacks the same witnesses
into the algebraicity, norm and torsion clauses of $\mathrm{IsKer}$.
See \path{FirstKernelLawReading.lean:74--79,127--130,428--434},
\path{PeriodLatticeFinitelyGenerated.lean:200--208,298--333} and
\path{GeneralLinearActions/IsDiscriminantTwelfthRootOfIdealLattice.lean:84--85}.

\paragraph{The coefficient identity after applying $\lambda$.}
\phantomsection\label{r14:comb-pseries}
$N_j=k+j(t_0+1)$ has even rank at both steps, since $k$ is even and the
rank parameter $t_0$ is odd. Both lattice objects in $t.O$ include their
self-duality witnesses. This gives the two evenness and two self-duality
clauses of admission. The reader's bad set is
$\{(D.\mathrm{mem}\,u).A.T,g_0,g_1\}\cup Y.Q_p(\mu)$.
The three clauses of \nolinkurl{LevFree} give $p\notin S$, and hence
$p\nmid\ell$. See \path{SatakeParameters/CombOfTied.lean:15--32},
\path{FirstKernelLawReading.lean:428--434}, \path{DataKerI.lean:65--67}
and \path{GeneralLinearActions/IsDiscriminantTwelfthRootOfIdealLattice.lean:171--174}.

For an admitted pair the unshifted $e$ is smooth to every complex
order. Indeed the standard vacuum belongs to its theta space by
$B.\mathrm{thZeroMem}$ under $B.\mathrm{StdOK}$. The two lattice moves
carry it into the theta space of $(L_0,L_1)$ by
$B.\mathrm{thMoved\_mem}$ and $B.\mathrm{LatAll}$
(p.~\pageref{r14:lat-indices}). Its first condition is complex
differentiability everywhere. Osgood's lemma in the coordinates of
\nolinkurl{analyticOnNhd_of_differentiable_vc}
(\path{AnalyticOnNhdOfDifferentiableVc.lean:21--32}) therefore supplies
all orders of smoothness, also after multiplication by the constant
character factor.
The curve ideal and its discriminant root supply $\delta\ne 0$ and
$B.\mathrm{StdOK}$. For every torsion shift, \nolinkurl{AlgTAll} at
$t.\mathrm{cc},t.\delta,t.B,t.\chi$ supplies algebraicity of every
Taylor coefficient of the translate. The character hypothesis here is
$t.\chi_{\rm ok}$. This gives precisely \nolinkurl{transcendence_algebraic_eigen}
and \nolinkurl{TransAlg}. See Sections~\ref{r7:Z4.2.04} and
\ref{r8:Z4.2.04}, \path{KernelLocalFactors/StdOK.lean:7--9},
\path{SelfDualLatticeNorms/SelfDualComap.lean:153--155},
\path{FirstKernelLawReading.lean:382--387,395--403,568--574} and
\path{PeriodLatticeFinitelyGenerated.lean:1826--1831}.

Write $a_\alpha(F)$ for the iterated derivative at zero in the repeated
coordinate directions of $\alpha$, divided by $\prod_j\alpha_j!$.
Each $U_v e$ is smooth by its translation and exponential formula.
Finite linearity of the iterated derivative, with the same factorial on
every term, gives $a_\alpha(\Theta)=\sum_{x,y}\rho_{x,y}
a_\alpha(U_{v_{x,y}}e)$. Every factor on the right belongs to
$\overline{\mathbf Q}$. Thus $\mathrm{lamC}$ agrees on the sum and on
each product with the ring homomorphism $\lambda$. For every $\alpha$,
\[
 [X^\alpha]\,\mathrm{pSer}_\lambda(\Theta)=
   \sum_{x,y}\lambda(\rho_{x,y})\,
   [X^\alpha]\,\mathrm{pSer}_\lambda(U_{v_{x,y}}e).
\]
This is formal coefficientwise equality, with no interchange of an
infinite sum. The witnesses in \nolinkurl{IsComb} take the \emph{same}
lattices $L_0,L_1$ at each summand. Both local agreement conditions
$\mathrm{EqAt}$ hold with multiplier $1$. Each $\lambda(\rho_{x,y})$
has norm at most one, and each translated point is killed modulo the
period lattice by the positive integer $\ell$ coprime to $p$.
This proves \nolinkurl{kernel_combination_form} on its stated
$\mathrm{TransAlg}$ and $\mathrm{IsKer}$ hypotheses (for general
$\mathrm{IsKer}$ use the possibly different $S$-supported annihilator
of each individual shift). Composing with the reader implication proves
\nolinkurl{first_kernel_law_reading} and specializes to
\nolinkurl{comb_of_tied} on the data above. See
\path{GeneralLinearActions/IsDiscriminantTwelfthRootOfIdealLattice.lean:115--118,178--181,211--213,280--295},
\path{KernelLocalFactors/ThDistC.lean:435--440} and
\path{PeriodLatticeFinitelyGenerated.lean:1457--1498,1520--1522}.

\paragraph{Use on the inert row and review boundary.}
The standard-basis integrality theorem chooses one unit $e_p$ and a
standard series $b_p$ at each qualifying place from the local host and
orbit hypotheses. It accepts any admitted pair, two standard bases and
an $\mathrm{IsComb}$ series there. For the two matrices of the adapted
reader, finite generation of each lattice, the adapted basis condition
and conjugate duality give the required standard bases. Transporting
$t.\mathrm{ent}.\mathrm{Th}$ by the adapted coefficient identity gives
the standard series of precisely the $p$-series above. Substitution gives
the required expansion for each orbit representative system $R$, with
sign function constantly $-1$. The choices of $e_p,b_p$ precede the class
pair and $R$; at a nonqualifying place the source takes $(1,0)$ in the
vacuous implication. See \nolinkurl{exists_beFam_evenConditionsInert_of_kerEnt},
\path{AdmissibleKernelSeriesInertPrimeNonConjugateFrob.lean:24--44,49--65}
and \nolinkurl{normalisedFamily_standardBases_integrality_alt},
\path{NormalisedFamilyStandardBasesIntegralityStatement.lean:12--13}.

Review B confirmed the conditional finite-expansion deduction in
Z4.3(ii). In addition to the inert-row caller above, three source
declarations use \nolinkurl{comb_of_tied}:
\nolinkurl{Ggeo.frobeniusConjugation_conditions_adaptedKernelDatum_alt}
(\path{FrobeniusConjugationConditionsAdaptedKernelDatum.lean:23--62}),
\nolinkurl{Ggeo.AdmEntInt.kernelDatum_entries_integral_away_from_d}
(\path{KernelDatumEntriesIntegralAwayFromD.lean:25--41}), and
\nolinkurl{Ggeo.KerShapeD.thetaKernelShapeAtD_ofKer_of_hyp_adapted}
(\path{KernelDatumThetaSeriesShapeAtDAdaptedBases.lean:16--42}).
All four pass the hypotheses $t.\mathrm{Hyp}$ and \nolinkurl{AlgTAll}.
In the latter two $\mathrm{AlgTAll}$ is the field $h_{\rm In}.\mathrm{algT}$.
The four outputs enter \nolinkurl{PdvClose.at_of_true_columns} through
their statement wrappers (\path{AtOfTrueColumns.lean:30--47}).
\nolinkurl{CurveThetaInert.torsionTaylorAlgebraic_latticeExists_of_cm}
supplies the translated Taylor law by \nolinkurl{ThLoc.algTAll},
and the lattice existence through CM2
(\path{TorsionTaylorAlgebraicLatticeExistsOfCm.lean:15--19};
\path{BundleBothPartsOfStandardKernelOfGlobalArtinReci.lean:15--24}).
The application assumes $t.\mathrm{Hyp}$; on the selected
true column the supplier takes $\mathrm{ColSide}(t)$, obtained from
$\mathrm{Duty}$, and $t.\mathrm{TopPinned}$. The adapted bases, the
local host and the split-bound hypotheses enter the application just
described. The dependencies on the actual selected row are recorded in
Section~\ref{r15:z43-premises}. The independent covariance argument
Z4.3(i) concerns $B.\mathrm{EqP}(\chi)$ at the inert and stabilizer
sites. The original Z4.3 input retains its accepted partial grade.

% Parent source SHA-256: 726efab7fd295a2dce4fc43b2a38c6ed087511bfb339c03e564ed0607ef7693c
% blueprint/ch_P_r15.tex lines 940 through 1003; r15:rational-weights
\paragraph{The two determinant exponents.}\label{r15:rational-weights}
Let $W$ be a Weil operator on $v$. Denote its two Hecke characters
by $\psi=W.\mathrm{chV}$ and $\psi'=W.\mathrm{chV}'$, their infinity weights
by $n_{\psi,\tau},n_{\psi',\tau}$, and the signs by $\epsilon,\epsilon'$.
Let $a_\tau,a'_\tau$ be integers satisfying
\[
 2a_\tau=\epsilon' n_{\psi',\tau}-m,\qquad
 2a'_\tau=\epsilon n_{\psi,\tau}+n .
\]
For rational unitaries $\gamma,\gamma'$ and their adelic images, every
all-level theta function $F$ satisfies
\[
 W(\gamma_A,\gamma'_A)F
 = \prod_\tau\bigl(\tau(\det\gamma)^{-a_\tau}
                      \tau(\det\gamma')^{-a'_\tau}\bigr)\,
 P_{\gamma,\gamma'}F,
 \qquad
 P_{\gamma,\gamma'}F(z)=F((\gamma'^*_\tau z_\tau\gamma_\tau)_\tau).
\]
Here the star denotes conjugate transpose after embedding.
See \nolinkurl{TaylorLogMetric.pull},
\path{GeneralLinearActions/IsDiscriminantTwelfthRootOfIdealLattice.lean:106--108},
and \nolinkurl{oddEq}, \path{GlobalSchurBounds/StandardLatticesGL3.lean:612--619}.

First act only by $\gamma$ with $1-\gamma$ nonsingular. Put $u=\det\gamma$.
Unitarity gives $u\bar u=1$. Write $u=d/\bar d$, taking $d=1+u$ for
$u\ne-1$ and a nonzero anti-invariant element for $u=-1$. At every
embedding, $(\tau(d)/|\tau(d)|)^2=\tau(u)$. The principal-idele law for
$\psi'$ and $2(a_\tau+m/2)=\epsilon'n_{\psi',\tau}$ then give the Weil
multiplier $\prod_\tau\tau(u)^{-(a_\tau+m/2)}$ on $v$ relative to the
normalized intertwiner. The analytic pullback on $v$ relative to the
same intertwiner has multiplier $\prod_\tau\tau(u)^{-m/2}$, from the
phase computation in Paragraph~\ref{r14:eq-rational}. The quotient gives
the first exponent. Acting only by $\gamma'$ with $1-\gamma'$ nonsingular
uses $\psi$ and $2(a'_\tau-n/2)=\epsilon n_{\psi,\tau}$. Its two
exponents on $\tau(\det\gamma')$ are $-(a'_\tau-n/2)$ and $+n/2$,
giving $-a'_\tau$ in the quotient. All multipliers divided out are nonzero.
See \nolinkurl{RatEqui.char_value}, \nolinkurl{RatEqui.det_root},
\path{FoundationProofs/OddEqOfEven.lean:26--71},
\nolinkurl{VacuumExact.vacuum_exact}, \path{WeilOperatorVacuumExactConstantEvenRank.lean:79--105},
\nolinkurl{RatRow.pull_eq_smul_iop}, \path{ThetaPullbackFirstGroupRationalPointEqSmul.lean:141--219},
and \nolinkurl{RatEqui.pull_right_eq_det_smul},
\path{ThetaPullbackSecondGroupRationalPointEqDetSmul.lean:150--234}.

Both operators intertwine the rational Heisenberg translations of $v$.
Their agreement there extends by linearity to all levels since these
translates span. To handle an arbitrary $\gamma$, factor it as
$(z\gamma)(\bar z\,1)$, where $z\bar z=1$, $z\ne1$ and $z^{-1}$ avoids
the finite spectrum of $\gamma$. There are infinitely many choices before
these exclusions, since $(r+\xi)/(r-\xi)$ are distinct as $r$ ranges
over $\mathbf Q^\times$ and $0\ne\xi=-\bar\xi$. Each factor has
nonsingular one minus the matrix. Multiplying the two identities gives
the formula for $\gamma$, and the same reasoning applies to $\gamma'$.
Finally the two actions compose by the group law and
$(\gamma'_1\gamma'_2)^*=\gamma_2'^*\gamma_1'^*$.
This gives all exponents in \nolinkurl{RatEqui.oddEq_of_even}.
Its $\mathrm{oddE0}$ supplies membership and $W.\mathrm{nz}$ supplies
nonvanishing. The two one-sided laws above therefore apply directly to any
nonzero standard theta vector even when only its membership is supplied.
See \nolinkurl{ThetaSpan.theta_span}, \path{ThetaSpanRationalTranslatesVacuum.lean:143--195},
\nolinkurl{RatExtend.om_eq_pull_of_vacuum}, \path{OmEqPullOfVacuum.lean:207--250},
and \nolinkurl{RatEqui.exists_factor}, \nolinkurl{IsPull.mul},
\nolinkurl{RatEqui.isPull_row}, \nolinkurl{RatEqui.isPull_col},
\path{FoundationProofs/OddEqOfEven.lean:82--109,117--180}.

% Parent source SHA-256: 2055d64ab882ca041a67ad1f4360db7dd6562c6c3f381c19b443acef9a2f98b1
% blueprint/ch_P_r9.tex lines 23 through 365; r9:Z4.3
\subsection*{Z4.3. Frobenius descent for the inert coordinate}
\phantomsection\label{r9:Z4.3}
At the inert-prime step of Section~\ref{r6:Z4.3}, put $K=\widehat{\mathbb Q_p^{\rm nr}}$,
$R=\mathcal O_K$, $q=p^2$, and $\tau=\sigma^2$, where $\sigma$ is the
isometric Frobenius lift. An exponent on $\tau$ acts on coefficients.
Congruence of two series means that each coefficient of their
\emph{difference} lies in $pR$. This convention applies even to
logarithms with unbounded coefficients. The formal-group calculation
below uses the following data.
\[
\begin{gathered}
 \lambda(0)=0,\quad\lambda'(0)=1,\quad
 \lambda(\Phi(X,Y))=\lambda(X)+\lambda(Y),\\
 \lambda(A)=a\lambda,\qquad \lambda(e)=p\lambda,\qquad
 |\sigma(a)-a|=1,\quad \tau(a)=a.
\end{gathered}
\]
Here $\Phi,A,e$ have integral coefficients and zero constant term,
and $e_i=[X^i]e$ is a unit for some $i\le q$.
The required estimate is for the specified $\lambda$ in this parameter.
Basenames in this calculation use the source directory
\path{ReciprocityTheory/ReductiveGroups/UnramifiedIntegrals/}. See
\path{LawD}, \path{HD}, \path{QCls} (\path{FgDefs.lean:14--46}) and
\path{FrobField}, \path{ResClosed} (\path{KurDefs.lean:22--37}).

\paragraph{Two coefficient tests.}
Write $V_\Phi=\partial_Y\Phi(X,0)$. The inverse of the normalized
logarithm gives $\Phi(X,0)=X$. Differentiating its addition formula
gives $\lambda'V_\Phi=1$. Since $V_\Phi\in R[[X]]$ and $V_\Phi(0)=1$,
$\lambda'$ is integral. If $f'$ is integral and two zero-constant
integral substitutions satisfy $P=Q+pB$, then $f(P)-f(Q)$ has
coefficients in $pR$. Indeed for $1\le j\le n$ the factor
\[
 f_n {n\choose j}p^j=(nf_n){n-1\choose j-1}\,p^j/j
\]
lies in $pR$, because $j-v_p(j)\ge1$. In each total degree the
binomial expansions of $f_n(Q+pB)^n$ are finite. This proof applies
in any finite number of variables.

Define $\delta D=\sum_{m\ge0}D_{pm}X^m$. If $D'$ has coefficients
in $pR$, then $D\equiv(\delta D)(X^p)$ and $(\delta D)'$ is
integral. Indeed, $nD_n\in pR$ forces $D_n\in pR$ for $p\nmid n$,
and $mD_{pm}=(pmD_{pm})/p\in R$. For an integral substitution
$P$ with zero constant,
$P(\mathbf X)^p\equiv(\sigma P)(\mathbf X^p)$, by Frobenius on
the residue field. Substitute $P$ in the congruence for $D$ and
apply the first test to $\delta D$. Thus
\[
 D(P(\mathbf X))\equiv(\delta D)\big((\sigma P)(\mathbf X^p)\big).
\]
Congruence after substitution of every variable by its $p$th power
reflects congruence, by reading the coefficients at the multiplied
multiindices. See \path{LawD.deriv_integral} (\path{FgHonB.lean:205--231}),
\path{qls} (\path{FgHonA.lean:137--192}) and \path{desc_congr},
\path{desc_key}, \path{desc_QCls} (\path{FgHonC.lean:44--148}).

\paragraph{Extracting the first unit and the successive twists.}
A series $D$ of zero constant and integral derivative belongs to
$Q(\Phi,A;b)$ if $D(\Phi)\equiv D(X)+D(Y)$ and $D(A)\equiv bD$.
Differentiation in $Y$ at zero gives
$D'V_\Phi\equiv D'(0)$, so $D'\equiv D'(0)\lambda'$.
If also $D'\equiv0$, the second test carries its two congruences
to $\delta D\in Q(\sigma\Phi,\sigma A;b)$ for $b\in R$.
For $D\in Q(\Phi,A;a)$ put $d=D'(0)$. Then $E=D-d\lambda$
is in the same class with $E'\equiv0$.
The linear term of $(\delta E)(\sigma A)\equiv a\delta E$
is $(\delta E)'(0)(\sigma a-a)\equiv0$, since $A'(0)=a$.
The unit difference kills this linear coefficient modulo $p$.
Applying the differential identity for $\sigma\Phi$ makes
$(\delta E)'\equiv0$. Descending again gives
\[
 D\equiv d\lambda+D_2(X^q),\qquad
 D_2\in Q(\tau\Phi,\tau A;a).
\]
This works at every $\tau$-twist, as $\tau a=a$ and $\sigma$ is
isometric. See \path{hon_deriv} (\path{FgHonB.lean:232}),
\path{eigen_trick}, \path{lemmaQ} (\path{FgHonD.lean:115--122,195--232})
and \path{lemmaQ_tw} (\path{FgHonF.lean:153--165}).

To initialize this iteration one needs a \emph{unit} tangent.
Applying the inverse logarithm shows that $e$ preserves $\Phi$ and
commutes with $A$. From $\lambda'(e)e'=p\lambda'$ we get
$e'\equiv0$, since $\lambda'(e)$ has integral inverse. Modulo $p$
write $e(X)=h_1(X^p)$ with $h_1=\delta e$.
The two exact homomorphism identities, reduced modulo $p$, say
$h_1(\sigma\Phi)\equiv\Phi(h_1(X),h_1(Y))$ and
$h_1(\sigma A)\equiv A(h_1)$. The latter identity gives
$h_1'(0)(\sigma a-a)\equiv0$. The former gives
$h_1'V_{\sigma\Phi}\equiv V_\Phi(h_1)h_1'(0)$, so $h_1'\equiv0$.
A second descent yields an integral $h=\delta^2 e$ with
$e\equiv h(X^q)$, $h(\tau\Phi)\equiv\Phi(h(X),h(Y))$ and
$h(\tau A)\equiv A(h)$. The unit of $e$ in degrees $1$ to $q$
must occur at $q$, so $h'(0)=e_q$ is a unit.
Our substitution test for $\lambda$ now gives
$D_0=\lambda(h)\in Q(\tau\Phi,\tau A;a)$,
$p\lambda\equiv D_0(X^q)$, and $D_0'(0)=e_q$.
Applying the two descents repeatedly produces $c_j\in R$,
$c_0=e_q\in R^\times$, with
\[
 p\lambda(X)\equiv\sum_{j\ge0}c_j(\tau^{j+1}\lambda)(X^{q^{j+1}}).
\]
In degree $k$, the unexpanded tail is a zero-constant series
in $X^{q^{j+1}}$ and eventually contributes zero.
See \path{LawD.subst_of_log}, \path{LawD.comm_of_log}, \path{hon_shape},
\path{hon_level} (\path{SubstOfLog.lean:15--56,88--135}),
\path{hon_mu} and \path{hon_series} (\path{FgHonF.lean:78--152,166--220}).

\paragraph{Eliminating the higher twists.}
Take $t_0=1$. There is a sequence $t_m\in R$ satisfying, for $n\ge2$,
\[
 p t_n=\sum_{i=0}^{n-1}t_i\,\tau^i(c_{n-1-i}).
\]
To construct it, send a bounded sequence $(s_m)_{m\ge1}$ in the unit
ball to the solution $(t_m)_{m\ge1}$ of the triangular equations
\[
 \sum_{i=1}^{m}\tau^i(c_{m-i})t_i=p s_{m+1}-c_m
 \qquad(m\ge1).
\]
The diagonal entries $\tau^m(c_0)$ are units. Forward substitution
keeps the solution integral. The same recursion on differences bounds
the sup norm by $|p|$ times the input difference. The unit ball of
bounded sequences is complete in the sup norm, so the contraction
has a fixed point, giving the desired recurrence.

For fixed $k>0$ put
$w_i=[X^k](\tau^i\lambda)(X^{q^i})$.
Twisting the preceding Frobenius relation and substituting $X^{q^i}$
gives $p w_i-\sum_{j\ge0}\tau^i(c_j) w_{i+j+1}\in pR$.
Multiply by $t_i$ and sum. Only $q^i\le k$ can contribute. The
recurrence cancels the coefficients of all $w_n$ with $n\ge2$.
What remains is $p w_0-(c_0-pt_1)w_1\in pR$.
Degree zero vanishes. Hence $c=c_0-pt_1$ is a unit and
\[
 \lambda(X)-(c/p)(\tau\lambda)(X^q)\in R[[X]].
\]
This proves the implication labelled \path{hon_main} in
Section~\ref{r6:Z4.3}. See \path{twisted_division}
(\path{FgHonG.lean:63--109}), \path{hon_regroup},
\path{hon_coeff_shift} and \path{hon_main} (\path{FgHonH.lean:15--77,192--256}).

\paragraph{Recovering the specified inverse coordinate.}
The residue field is $\overline{\mathbf F}_p$. Lift a nonzero root
of $z^{q-1}=\bar c^{-1}$ to a unit $u$. If $u-c\tau u=p^m r$
for $m\ge1$, replace $u$ by $u+p^m v$ with
$\bar v-\bar c\bar v^q=-\bar r$. Such a root exists in the same
residue field, and the new error lies in $p^{m+1}R$. Completeness
and continuity give a unit $\epsilon=c\tau\epsilon$.
Set $f=\epsilon^{-1}\lambda$ and
$L(X)=\sum_{j\ge0}X^{q^j}/p^j$. Then
$B=f-p^{-1}(\tau f)(X^q)$ is integral and $L-p^{-1}L(X^q)=X$.
To construct $H$ with $f(H)=L$, start with $H=0$, and suppose
an integral polynomial $H$ of zero constant has matched the terms
below degree $n\ge1$. We have
\[
 L-f(H)=X-B(H)+p^{-1}\big(L(X^q)-(\tau f)(H^q)\big).
\]
The coefficient of degree $n$ on the right is integral. Indeed
$H^q\equiv(\tau H)(X^q)$, and $(\tau f)'$ is integral. Replacing
$H^q$ by this congruent series changes the parenthesis by a series
in $pR[[X]]$. The resulting parenthesis is
$\tau(L-f(H))(X^q)$ and vanishes through degree $n$, since $qn>n$.
One may therefore add $[X^n](L-f(H))/f'(0)$ times $X^n$ to $H$,
as $f'(0)=\epsilon^{-1}$ is a unit. Induction proves
$H=\lambda^{-1}(\epsilon L)\in R[[X]]$ with linear coefficient
$\epsilon$. Under $K\hookrightarrow\mathbf C_p$ this is precisely
$u_{g_2,g_3}(\epsilon L)$, since the short exponential and logarithm
are compositional inverses. The choice of $\epsilon$ precedes the
coefficient index, as required by CM3. See
\path{lang_semilinear} (\path{SemilinNormPowSubPow.lean:76--160}),
\path{fin_integral} (\path{FgFinA.lean:115--178}) and
\path{cm3} (\path{FgClose.lean:25--91}).

\paragraph{Checking the height-two hypotheses on the specified short model.}
Fix squarefree $d\equiv3\pmod4$, the CM period lattice, an inert prime $p\geq5$ with $p\nmid d$,
the specified $\iota:\mathbf C\simeq\mathbf C_p$, and \emph{any}
good model $(g_2,g_3)$ for them. Write $U$ for the completion of the
maximal unramified extension of $\mathbf Q_p$ inside $\mathbf C_p$.
Its residue field is $k=\overline{\mathbf F}_p$. Indeed the nonzero
elements of each $\mathbf F_{p^s}$ lift uniquely to roots of
$X^{p^s-1}=1$ in its unramified lift. These roots, for all $s$,
generate the dense unramified field. Completion preserves the residue
field and value group because a Cauchy sequence with nonzero limit has
eventually constant norm, and can be approximated modulo $p$.
Arithmetic Frobenius on this field sends a prime-to-$p$ root $\zeta$
to $\zeta^p$. It is isometric and extends continuously to $U$, giving
$\sigma(x)\equiv x^p\pmod p$ on its valuation ring $R$.
For $a=\iota(\omega_d)$, the equation $a^2-a+(d+1)/4=0$ reduces to
an irreducible separable quadratic since its discriminant is $-d$.
Both roots have unique Hensel lifts in $R$. Frobenius interchanges
them, so $\sigma a=1-a$, $|\sigma a-a|=1$ and $\sigma^2 a=a$.
See \nolinkurl{kur_complete} and \nolinkurl{FrobField}
(\path{KurDefs.lean:14--36}), \nolinkurl{kurFrob}
(\path{KmuInclDenseRange.lean:26--67}), \nolinkurl{kur_resClosed}
(\path{TchOfIntegral.lean:25}) and \nolinkurl{kur_quad}
(\path{NormSqrtNegPowPrimeSqSub.lean:177}).

The defining scale gives a homothetic complex lattice with invariants
mapping to $g_2,g_3$ and with the same endomorphism $z\mapsto\omega_d z$.
The equation $y^2=x^3-g_2x/4-g_3/4$ over $R$ has unit discriminant.
It defines a smooth proper elliptic scheme $E$. Its group law completed
at the origin is integral, and in the short-curve parameter its
normalized logarithm and inverse are $\lambda=\mathrm{lgS}$ and
$u=\mathrm{uSer}$. This follows by integrating the invariant
differential normalized to have linear coefficient one, or by inverting
the uniformization parameter $-2\wp/\wp'$. \emph{Alternative proof of endomorphism integrality.} The algebraic endomorphism
of the uniformized complex curve transfers under $\iota$ to one with
tangent multiplier $a$. To see it over a finite extension of $U$,
put $m=(d+1)/4$. Its kernel lies in $E[m]$ since
$\omega_d(1-\omega_d)=m$. A finite extension defines its kernel
points and hence the quotient isogeny. The induced isomorphism
from the quotient back to $E$ also descends after finite extension:
in characteristic zero put both curves in short form, where an
origin-preserving isomorphism is given by a scaling whose fourth or
sixth power is prescribed by a nonzero coefficient. Over the new
valuation ring both source and target of the endomorphism have good
reduction. The N\'eron mapping property extends it, and uniqueness
preserves the group law and origin. Its formal series is
$A=u(a\lambda)$ by the logarithm and the tangent multiplier. Each
coefficient belongs to $U$ by this formula and is integral by extension,
so $A\in R[[T]]$. The mapping property used here is stated in
Conrad \cite[Definition~5.1 and Example~5.3]{P-R9-Conrad}. The source uses integral coordinates of the nonzero
kernel points to reach the same conclusion. See \nolinkurl{good_lattice}
(\path{FgAnLat.lean:11--27}), \nolinkurl{lgS_eq_logSer} and
\nolinkurl{curve_law} (\path{ShortWIntModel.lean:12--35}),
\nolinkurl{endo_curve} (\path{FgAnEndo.lean:11--48}) and
\nolinkurl{end_integral} (\path{EndIntegral.lean:54--115}).

Iterate the integral formal addition law, starting from $[0]=0$ and
$[n+1]=\Phi([n],X)$. Each $[n]$ is integral and
$\lambda([n])=n\lambda$ by induction. Inverting the logarithm gives
$e=[p]_E=u(p\lambda)$ (\nolinkurl{mulN_integral}, \path{FgMulN.lean:26--88}).
It reduces to multiplication
by $p$ in the formal group of $E_k$.

\emph{Alternative proof of the degree bound.} The morphism $[p]$ of this
elliptic curve is finite flat of degree $p^2$ \cite[Lemma~39.9.8, Tag~0BFG]{P-R9-Stacks}. Its full kernel therefore has length $p^2$.
The completed local algebra of that kernel at the identity is
$k[[T]]/(\overline e(T))$, so $\overline e$ has a first nonzero
coefficient of degree at most $p^2$. Its lift is a unit in $R$.
The Lean proof instead uses \nolinkurl{good_lattice}, \nolinkurl{half_cubic}
and \nolinkurl{half_prod_norm} in \nolinkurl{level_le} (\path{FgLevel.lean:41--87})
used in \nolinkurl{cm3} (\path{FgClose.lean:25--91}),
before determining the formal height.

\paragraph{The theta identity still required on the selected row.}\phantomsection\label{r9:z43-open}
Fix $D$, the curve ideal $\mathfrak c$, its discriminant twelfth root
$\delta$, the base $B$ on this datum and a character $\chi$ satisfying
$\mathrm{IsChi}(D,\chi)$. Write $S,S'$ for the two standard lattices
of even ranks $N,N'$. At an inert $p\ge5$, weak approximation expresses
an admitted pair as $(\gamma M,\gamma'M')$, with $M,M'$ self-dual and
equal locally at $p$ to $S,S'$, and $(\gamma,\gamma')\in
U_N(E)\times U_{N'}(E)$ rational. For each such choice used in the
reduction, with $\gamma_\tau,\gamma'_\tau$ denoting the embedded
matrices, the following equality of entire functions is needed:
\[
 e^{\rm nat}_{B,\chi}(\gamma M,\gamma'M')(z)=
 e^{\rm nat}_{B,\chi}(M,M')
       \bigl((\gamma_\tau^{\prime *}z_\tau\gamma_\tau)_\tau\bigr).
\]
There is no scalar in this change of variables. Explicitly, $e^{\rm nat}$
is the standard tensor theta acted on by the two normalized lattice
averages along $(S,S')\to(M,S')\to(M,M')$, and multiplied by
\[
 C_\chi(M,M')=
 \chi(\mathrm{detId}(S,M))^{N'/2}
 \chi(\mathrm{detId}(S',M'))^{-N/2}.
\]
Each average has the inverse square root of the index and the
semicharacter times the translated vacuum in its sum. Reindexing sends
the starting pair to $(\gamma S,\gamma'S')$. To return to $(S,S')$
requires its vacuum scalar and the compatibility of the intervening
averages. For instance the one-sided standard-vacuum factors in the
direct proof are
$(\prod_\tau\tau(\overline{\det\gamma}))^{N'/2}$ on the source
and $(\prod_\tau\tau(\det\gamma'))^{N/2}$ on the target. The rule
for $\chi(\mathrm{detId}(S,\gamma M))$ multiplies its old value by
$\prod_\tau\tau(\det\gamma)^{-1}$. Since a unitary determinant has
$\overline{\det\gamma}=(\det\gamma)^{-1}$, the ratio
$C_\chi(\gamma M,\gamma'M')/C_\chi(M,M')$ is the product of the two displayed factors.
The normalization cancels only if the inverse of that product occurs when the
unnormalized vacuum is expressed after moving both reference lattices back.
Verifying this return map, including its scalar on the generators, remains open
in the hand derivation.
See \path{KernelLattices/LateralTheory.lean:165--178},
\path{GeneralLinearActions/IsDiscriminantTwelfthRootOfIdealLattice.lean:100--109},
\path{GlobalSchurBounds/TaylorLogMetric.lean:37--50},
\path{SplitPairEquations/Reduction.lean:74--94} and
\path{NormalisedFamilySatisfiesEquivarianceLaws.lean:17--31}.

Given the displayed equality, transport of a torsion point and a
local basis, integrality at the standard pair and finite linearity
give the desired estimate as in Paragraph~\ref{par:r4-z43-host}.
The source invokes the first conjunct of
\nolinkurl{eulerProduct_eq_of_curve_ideal} precisely for this equality
(\path{InertCoprimeLemmas/IntAsm.lean:55--77},
\path{SplitPairEquations/FinalEq.lean:6--9}). The direct proof uses
the square and triangle identities and the two standard-vacuum
scalars; its calculation on generators passes through the finite Gauss
sum (\path{IntegralSquareRoots/SecondConstantOfTrianglesData.lean:14--32},
\path{IntegralSquareRoots/SecondConstantOfGaussSum.lean:14--31}).
The finite-sum argument on the fixed target stabilizer already used in
this chapter assumes this very covariance for $(1,\gamma')$
(\path{SatakeParameters/ThCLeftStab.lean:368--372,435--436}).
The alternative reference-family deduction on
page~\pageref{r8:Z4.2.13} requires constructing a family with its three
move laws and comparing its standard vector, as seen in
\path{NormalisedFamilySatisfiesEquivarianceLaws.lean:59--99}.
The normalization for the rational moves above is therefore
\emph{[Proved in Lean, not reproduced]}.
This is an internal call of the inert integrality supplied to the
row from CM2 (\path{IntInertOfComplexMultiplicationCondition.lean:11--14},
\path{BundleBothPartsOfStandardKernelOfGlobalArtinReci.lean:15--24}).

\paragraph{What the target stabilizer fixes.}
For a word, column, step $i<m$ and class pair $(c,c')$ of the auxiliary true-column reader on the same chain, put $N=r.N_i(c)$, $N'=r.N_{i+1}(c')$. Write $e_{N,N'}$ for its seam-reindexed normalized vacuum. The target stabilizer is represented by a rational unitary $\gamma$ with $\gamma N'=N'$. Its action $L_\gamma$ on the archimedean variables is left multiplication, block by block, by the embedded matrix of $\gamma$. The stabilizer calculation requires
\[
             e_{N,N'}(L_\gamma z)=e_{N,N'}(z)\qquad(\text{every }z).
\]
Indeed the two lattice arguments in simultaneous unitary equivariance become $N$ and $N'$ on taking the source matrix to be $1$. Transport across the seam identifies the inverse pullback matrix with the conjugate transpose of the embedded target matrix. This gives the displayed equality by specializing the same first component of $B.\mathrm{EqP}(\chi)$ used for the inert transfer. Once this equality holds, reindexing the kernel under $L_\gamma$ leaves only the finite-datum multiplier $\sigma_0(k)/M$. The implementation supplies equivariance by comparing the normalized vacuum to a reference family and invoking its simultaneous change-of-basis law. The restricted invariance for these pairs is \emph{[Proved in Lean, not reproduced]} as described at Z4.2.13.
See \path{SatakeParameters/PairKernelLawOfFiniteDatumLaw.lean:54--61,85--104}, \path{SatakeParameters/ThCLeftStab.lean:368--397,432--436} and \path{IwasawaDecomposition/NormalisedFamilySatisfiesEquivarianceLaws.lean:59--76,89--101}.

\paragraph{Which finite sums enter the row.}
For coefficient-field rationality the standard reader writes the supported periodic datum as a sum of its values times coset indicators, one for each class of the divided Hom lattice modulo the multiplied Hom lattice. Linearity gives the sum of indicator kernels once each induced indicator kernel is identified at the product level $c_Ac_B$. Only after this identification does the supplied \path{JetsIn} place its normalized jets in the field. See \path{SatakeParameters/TrueKernelJetsMulDetCharactersMemCoefficientFiel.lean:28--55,57--68}, \path{SatakeParameters/ThetaSupportInDivLattice.lean:81--110} and \path{KernelLocalFactors/ThDistC.lean:1768--1769}.

For the adapted, level-free series of the actual tied reader, \path{comb_of_tied} uses a second identity. Each selected kernel must expand into translates of $B.eNat(\chi)$ with algebraic weights integral away from the bad set and shift orders supported there. The source inserts three sets of representatives, normalizes the middle Fourier sum and combines the outer shifts by the Heisenberg phase. At volume one, write $\Phi$ for the selected finite datum, $\alpha$ for the semicharacter, $E_H$ for the alternating form and $S_X,S_A,S_B$ for the three representative sets. The required terms are $U_{x+y}e$ for $x\in S_X$, $y\in S_B$, with weights
\[
 r_{x,y}=\frac{\alpha(x)\alpha(y)}{|S_A|}
    e^{-\pi i E_H(x,y)}
    \sum_{t\in S_A}\Phi(t-x)\,e^{2\pi i E_H(t,y)}.
\]
This follows formally by substituting \path{hatS} and applying \path{opU_xiS} (\path{KernelLattices/XiIndCClass.lean:18--28}, \path{KernelLocalFactors/PeriodLatticeFinitelyGenerated.lean:168--196}). The outstanding specialization must identify these representatives with the three quotients of the selected lattice pair, prove that every $r_{x,y}$ is algebraic and integral outside its bad set, and show that $x+y$ is killed by an integer supported on that set.
The phase and the normalization for the selected adjacencies are precisely the specialization of the calculation left open at Z4.2.04. This use is \emph{[Proved in Lean, not reproduced]}. It enters the inert bound through \path{SatakeParameters/AdmissibleKernelSeriesInertPrimeNonConjugateFrob.lean:59}, \path{SatakeParameters/CombOfTied.lean:14--33} and \path{KernelLocalFactors/PeriodLatticeFinitelyGenerated.lean:179--196,276--290,320--333,1520--1522}.

\paragraph{Verdict for the original input.}
The Frobenius comparison and the integral coordinate now hold for the specified good model.
The two open arguments for the actual column are the first conjunct of
$B.\mathrm{EqP}(\chi)$, used at both the inert and the stabilizer sites, and the
Z4.2.04 specialization of the weighted kernel expansion in \nolinkurl{comb_of_tied}.
The direct covariance and finite expansion have cones of 24,010 and 15,441
core lines. Even conditional on both arguments, the application retains the
the application premises registered on p.~\pageref{r10:z43-register}.
Z4.3 remains \emph{partial}. Thirteen of the twenty original inputs meet the
conditional completion criterion in the trust ledger.

